\chapter{Discussion}

The empirical findings of Chapter 4 admit several interpretations, each with implications for how RAG evaluation should be designed and how its trustworthiness should be measured. This chapter discusses four threads in turn. Section 5.1 unpacks the central paradox finding and what it means mechanistically. Section 5.2 examines the cross-judge differences and what they imply for the LLM-as-a-Judge paradigm. Section 5.3 places the work in the context of existing RAG defences. Section 5.4 outlines mitigation directions the present results point toward. Section 5.5 records limitations. Section 5.6 considers practical implications for practitioners deploying LLM-as-a-Judge pipelines.

\section{The Filtering Paradox: Mechanism and Interpretation}

The headline empirical finding is that pre-filtering does not reliably mitigate judge vulnerability to corpus poisoning, and on the most adversarial attack class it substantially shifts judge behaviour in a counterintuitive direction. The largest and most statistically robust effect is the DeepSeek Survived-subset shift on PoisonedRAG-style triplets, where mean faithfulness rises from 0.361 at baseline to 0.738 after the statistical filter (Δ +0.377; question-clustered bootstrap p = 0.012). The True-subset effects are smaller and less robust statistically. For DeepSeek, the True Δ of +0.151 is borderline-significant after question-level bootstrap correction (p = 0.099). For GPT and Gemma, the True Δs of +0.053 and +0.072 do not survive clustered correction (bootstrap p = 0.239 and 0.259 respectively). The directional pattern - filter raises True faithfulness for all three judges - is real, but only the largest effects survive proper statistical scrutiny.

This finding is counterintuitive on its face. If a filter is a strong classifier (the statistical filter achieves F1 = 0.929 on held-out passage-level classification), removing the poison it identifies should make the judge's job easier, not harder. Section 4.5 tested two candidate mechanisms empirically and surfaced a third. This section unpacks what we found.

\subsection{Collateral damage is real but it is not the operative mechanism}

Collateral damage - the filter removing supporting passages it should have spared - happens at a meaningful rate (roughly 17.5\% of supporting passages across the dataset). The natural intuition is that this is what drives the True-subset paradox: the filter takes away the evidence the judge would have used to detect the contradiction, and so the surviving poison goes uncontested.

Section 4.5.2 tested this directly. Within the True subset, we partitioned triplets into those where the filter removed at least one supporting passage and those where it removed zero. If collateral damage drives the paradox, the support-removed group should show the largest positive Δ. The result is the opposite. In all three judges, when the filter removed supporting passages, faithfulness scores fell by 0.15 to 0.23. The judge correctly registered the missing evidence. The paradox - post-filter scores rising - emerges only in triplets where the filter spared the supporting passages.

So collateral damage exists as a phenomenon, but it operates in the opposite direction to the paradox. When it happens, the judge handles it the way an unbiased reasoner would: less evidence, less confidence, lower score. The aggregate True-subset Δ reported in Chapter 4 is the net of two opposing effects: a small negative pull from the 27\% of triplets where collateral occurred, and a larger positive push from the 73\% of triplets where the supporting passages survived the filter.

\subsection{The operative mechanism is context clarification through distractor removal}

What actually drives the paradox is a different effect entirely. The statistical filter, designed to identify poisoned passages, also removes a substantial fraction of the original HotpotQA distractor passages - passages that the dataset constructors deliberately included to make the multi-hop task harder. These distractors are noisy, topically loose, and embedding-distant from the question. They share surface features with the random\_noise injection type, and the filter learns to remove them along with the actual poison.

The downstream effect on the judge is straightforward in retrospect. The post-filter context is shorter (mean 7.1 passages versus 11.6 at baseline, a 39\% reduction). It is more topically focused, because the filter removed the most off-topic passages first. The remaining context - including any malicious passages the filter did not catch - sits inside a presentation that reads as cleaner, more coherent, and more answer-relevant. The judge reads this cleaner residual and assigns it a higher faithfulness score.

Section 4.5.2 reports the stratification: when we partition True-subset triplets by the quartile of original distractors removed, the Δ rises monotonically with distractor removal in all three judges. In the lowest quartile (few distractors removed), there is no paradox. In the highest quartile (many distractors removed), the Δ reaches +0.25 to +0.46. The relationship is not a small statistical artefact; it is the dominant pattern in the data.

This corresponds to a length effect that is independently visible in the judge's scoring behaviour. The Pearson correlation between context length and faithfulness score is weak at baseline (r = 0.10 for GPT, 0.29 for Gemma, −0.08 for DeepSeek) and substantially stronger post-filter (r = 0.36, 0.47, 0.31). After the filter has run, judges score shorter contexts as more faithful - a behavioural pattern that is not present, or barely present, before filtering.

\subsection{Why this is a paradox in the practical sense}

We continue to call this a paradox because the local effect of filtering (removing identified poisoned passages) and the global effect (judge reliability) move in opposite directions for a meaningful fraction of cases. But the mechanism is now sharper than the literature on filter-then-evaluate pipelines has assumed.

The standard intuition runs: poison degrades the residual, the filter removes poison, so the residual improves and the judge gets a better signal. This thesis shows that the residual produced by filtering is improved in a way that is hostile to careful judgment - it is shorter, cleaner-looking, more topically focused, and therefore more credible than the original noisy context, regardless of whether the surviving passages are actually faithful to the answer. The judge trusts the residual because it reads as well-formed, not because it actually contains the supporting evidence.

This reframes the engineering problem. Filter design has been optimised against passage-level classification metrics (precision, recall, F1) on the assumption that better classification produces better downstream judgment. The data here shows that the link from classification quality to downstream judgment is not monotonic. A filter that aggressively removes "noisy-looking" passages will produce residuals that judges find very convincing - including the ones that contain residual poison - which is exactly the wrong direction.

\subsection{Implications for filter design}

Three implications follow.

First, filter design has been optimising the wrong objective. Classifier-style metrics measure how well the filter identifies poisoned passages. They do not measure how the filtered residual interacts with downstream judges, and the data shows that those two things can diverge substantially. The right metric is end-to-end judge response on a held-out test set with realistic poisoning rates, not classifier F1 on a held-out passage-level test set.

Second, the distinction between "removing poison" and "removing distractors" matters in deployment. Filters that aggressively prune off-topic content will catch some poison but will also strip out original retrieved noise - context that, paradoxically, may have been useful precisely because it gave the judge something to ground a sceptical reading in. Filters designed to preserve original context shape (re-ranking rather than removal, contradiction detection rather than per-passage classification) face different trade-offs that may be more favourable to downstream judgment.

Third, the dependency of the paradox on context length suggests that judges have a length-based credibility heuristic that filtering activates. This is not a mistake of the judges per se - shorter, more focused contexts are usually more trustworthy in non-adversarial settings - but it is exploitable when the residual is the product of a filter rather than of an honest retrieval. Judges aware of having been pre-filtered might calibrate differently, and Section 4.6's small qualitative justification analysis offers a hint that prompt-level interventions can shift this. Quantifying the effect at scale is future work.

\section{What the Cross-Judge Differences Mean}

The three judges studied here are roughly contemporary and span the proprietary, open-weight, and cost-efficient segments of the LLM market. They differ in size, training data, and architectural choices, but all three are competent at standard NLP tasks. The differences observed in Chapter 4 are not differences of basic capability.

The differences are differences of scoring strategy. GPT's continuous calibration reflects a graded interpretation of the [0, 1] interval: the judge actively modulates its confidence as a function of evidence strength. Gemma's extreme leniency reflects a default-faithful prior: scores cluster near 1.0 unless evidence strongly contradicts the answer. DeepSeek's near-binary calibration reflects a categorical decision pattern: the judge has decided that the answer is faithful or it has decided that it is not, and the score reports the verdict at one of the extremes.

These three regimes are not equally good or equally bad. Each carries trade-offs.

GPT's continuous scoring is most useful for downstream pipeline design because it preserves the full interval as a confidence signal. A pipeline can route triplets with scores in [0.4, 0.7] to human review while passing triplets near the extremes. Gemma's leniency means high faithfulness scores are uninformative - they may indicate genuine faithfulness or they may indicate the default. Low scores from Gemma are valuable signal but rare. DeepSeek's near-binary scoring is the inverse problem: every score is a confident verdict, but the confidence is overconfident - the judge is not actually as certain as the score suggests, as evidenced by the substantial baseline FPR that accompanies the binary scoring style.

The cross-judge directional reversal on PoisonedRAG-style attacks is the most striking single observation from Chapter 4. GPT and Gemma find this attack hardest; DeepSeek finds it easiest. The reversal shows that the relationship between attack sophistication and detection difficulty is not monotonic across judges. An attack designed to be subtle (PoisonedRAG-style is engineered to look topically coherent and to embed retrieval anchors) can be the easiest to detect if the judge's scoring strategy happens to respond to a specific structural feature of the attack.

DeepSeek's near-binary scoring effectively functions as a contradiction detector on PoisonedRAG-style triplets. The malicious passage explicitly supports an answer that contradicts the ground-truth answer in the answer slot. Once DeepSeek detects the contradiction, its scoring style commits to a low score. Continuous-scoring judges weigh the contradiction against the topical fluency of the malicious passage and arrive at a high score. Same evidence, different verdicts.

Two implications follow.

Judge selection is a security-relevant decision. Most published evaluation work selects judges based on cost or availability and treats them as broadly interchangeable for the same task. The results here suggest, on a sample of three judges and a single retrieval-QA task, that a deployed pipeline's robustness to corpus poisoning can depend substantively on which judge it uses, and that the "right" judge depends on the threat model. We do not claim to have shown this for the LLM-as-a-Judge paradigm at large - three judges and one task are too narrow a base for that - but the directional point is concrete enough that practitioners should measure baseline robustness under representative adversarial conditions before committing to a judge for production use.

The literature's standard practice of reporting average results across judges may mask important per-judge differences. Section 4.1.3 would have looked very different if presented as a "mean baseline FPR across three judges": the directional reversal would have been hidden. Multi-judge studies should report per-judge results explicitly and should resist the temptation to summarise across calibration regimes that respond to attacks differently in kind.

We do not argue that one calibration regime is uniformly preferable to the others. Each has trade-offs, and the right choice depends on the threat model and the pipeline's downstream design. The practical recommendation is that pipelines deploying LLM-as-a-Judge should characterise the calibration of their chosen judge under representative adversarial conditions before deployment, and should account for that calibration when designing downstream routing or aggregation logic.

\section{Comparison with Existing RAG Defences}

The RAG security literature contains several published defence mechanisms, all of them aimed at the generator. We briefly survey the main families and contrast them with the present work.

Detection-based defences (perplexity, embedding clustering) assume that adversarially crafted text exhibits unusual statistical properties. The original PoisonedRAG paper showed that the perplexity distributions of clean and malicious texts overlap heavily, and the PR-Attack framework deliberately minimises perplexity, rendering this defence essentially inert. Embedding clustering defences (e.g., TrustRAG) assume that malicious documents cluster distinctly in embedding space; this works at low poisoning rates but degrades when attackers outnumber benign passages. Both families operate at a similar abstraction level to one signal of our statistical filter, and our filter inherits their strengths and limitations on that dimension while compensating with other signals.

Aggregation-based defences (RobustRAG) run the generator independently on each retrieved passage and aggregate the results via voting. The defence is effective when poisoned passages are a minority of retrievals, but fails when attackers can guarantee retrieval of multiple poisoned passages. Aggregation operates at the generator side; it does not address judge reliability.

Adversarial training (RAGuard) trains the retriever or generator to be robust to adversarial documents. The approach is elegant but requires representative attack data, which limits transferability to novel attack types. It does not transfer to the judge.

Process-based defences (InstructRAG, AstuteRAG, Discern-and-Answer) modify the generation process itself, asking the LLM to reason about reliability or to explicitly reject poisoned context. They are effective against denial-of-service attacks but less effective against targeted poisoning.

Lightweight ML post-retrieval (RAGDefender) applies a lightweight ML filter to retrieved documents before they reach the generator. It limits Attack Success Rate to single-digit percentages on PoisonedRAG-style attacks. The architecture is closer to ours than the others - a lightweight filter between retrieval and downstream consumption - but the consumer in their work is the generator and the metric is ASR.

The pre-filter studied in this thesis is closest in spirit to RAGDefender. The differences are: our consumer is the LLM-as-a-Judge rather than the generator; our metric is judge FPR across three RAGAS dimensions rather than ASR; we evaluate two filter architectures (statistical and LLM-based) for cross-architecture comparison; and we use the True/Survived decomposition to measure judge reliability under operationally correct labelling.

The empirical finding distinguishing this work is that, when measured against judges rather than generators, the same filter family that succeeds for generator defence does not reliably help. This is not a refutation of the existing defence literature but an extension into a setting that has not previously been measured. Defending the judge is a distinct problem, and the present work establishes that it is not solved by the current generator-focused defence toolbox.

\section{Mitigation Strategies the Results Point Toward}

The empirical results suggest several directions for further mitigation work. We outline four here, in roughly increasing order of effort to implement.

\textbf{Hybrid filter architectures.} The statistical filter and the LLM filter detect partially complementary phenomena. The statistical filter has near-perfect recall on PoisonedRAG-style and somewhat lower recall on adversarial fact. The LLM filter has the opposite profile - strong on PoisonedRAG-style and adversarial fact, weak on random noise. A hybrid that uses the statistical filter as a first-pass sieve and routes ambiguous triplets to the LLM filter would, in principle, combine the strengths of both. Whether the gain in classification F1 translates to improved end-to-end judge reliability is the open question. The mechanism analysis in Section 5.1 suggests the answer is not automatic: a hybrid that achieves higher recall may produce shorter, cleaner-looking residuals through more aggressive distractor removal, recreating the paradox at a higher operating point.

\textbf{Calibration-aware judge selection and ensembling.} Section 5.2 noted that the three judges have qualitatively different calibration regimes and respond differently to attack types. A practitioner deploying an LLM-as-a-Judge pipeline could explicitly select for calibration properties suited to the threat model. For threat models dominated by PoisonedRAG-style attacks, a near-binary judge with strong contradiction detection (analogous to DeepSeek) is the better choice. For threat models dominated by random noise or factual-edit attacks, a continuous-scoring judge with confidence-abstention support (analogous to GPT) is preferable. A pipeline could even ensemble different calibration regimes - not for variance-based detection (which Section 4.7 ruled out) but for coverage of different failure modes. The cost is computational: running three judges per triplet rather than one.

\textbf{Confidence abstention for continuous-scoring judges.} GPT's continuous calibration creates a natural decision band. A pipeline could abstain on triplets with faithfulness scores in (e.g.) [0.4, 0.7] and route them to human review. Score distributions in Section 4.1.2 suggest a 10–15\% abstention rate would capture most of the residual judge errors after filtering. The cost is human review load on the abstained subset, but the benefit is that abstention is exactly where post-filter judge errors are concentrated. We did not implement and measure this strategy in the present work, but it is a tractable empirical question for follow-up.

\textbf{Judge-side hardening.} If pre-filtering cannot reliably produce residuals that judges handle well, the alternative is to harden the judge directly. Possible directions include: training judges on adversarial examples so they learn to recognise polished poison; modifying the scoring prompt to require explicit cross-passage consistency reasoning before scoring; adding a meta-evaluation layer that scores not just the answer but the trustworthiness of the context. The work in Section 4.6 (justification analysis) shows that explicit poison-aware prompting shifts judge behaviour substantially on PoisonedRAG-style attacks. This is the most promising single direction the data supports, and it does not require any change to filter design.

These directions are not mutually exclusive. A production pipeline could combine all four: a statistical filter as the first pass, a triplet-level LLM filter for ambiguous cases, confidence abstention on continuous judges, and judge-side prompting that requires explicit consistency reasoning. The empirical question is which combinations achieve the best cost-benefit ratio. That is an evaluation programme of its own.

\section{Limitations}

The study has several methodological limitations worth recording explicitly.

\textbf{Single-domain evaluation.} All experiments use HotpotQA, and the findings should be read as "how three LLM judges behave on multi-hop retrieval QA under controlled poisoning" rather than as generic claims about LLM-as-a-Judge. The attack types are tuned to a multi-passage retrieval setting where supporting evidence is distributed across multiple cooperating passages. The pre-filter signals (cosine anomaly, answer-span recall, cross-encoder relevance) exploit properties of this setting. Whether the headline paradox - filtering raises True-category faithfulness - generalises to single-hop QA, dialogue evaluation, summarisation faithfulness scoring, or open-domain agentic settings is genuinely unknown. The closest reasonable extrapolation is to other multi-passage retrieval tasks; extrapolation beyond that is speculation. We treat single-task evaluation as the most consequential scope limit on the conclusions drawn here.

\textbf{Three judges, not a representative panel.} Our judge sample is three contemporary LLMs from three different providers (one proprietary, one open-weight, one cost-efficient). With n=3 we cannot make distributional claims about how LLM judges as a class behave. We can demonstrate that \textit{these} three judges occupy distinct calibration regimes, that \textit{these} three respond differently to the most sophisticated attack class, and that \textit{these} three exhibit the True-subset paradox under filtering. Whether a fourth or fifth judge - a reasoning model, a frontier-tier model, a fine-tuned domain-specific judge - would extend, weaken, or break the patterns is an empirical question we cannot answer from our data. Statements throughout Chapter 5 such as "judge architecture matters as much as judge size" should be read as suggested by the panel, not established by it. A larger panel that includes reasoning models is the most direct extension of this work.

\textbf{No human-annotated ground truth.} The thesis evaluates LLM judges against each other and against the construction-time poisoning labels. It does not validate the judges against human assessment of faithfulness on the same triplets. The standard RAGAS premise that LLM judges proxy human judgement is therefore inherited from the literature rather than verified in this work. Our findings are best read as "how three LLM judges' verdicts shift under attack" rather than "how judgement quality, in any human-validated sense, changes under attack." A small human-annotated subset (say 50 triplets, 3 annotators) would let future work make the stronger claim. We treat this as the second most consequential scope limit on the conclusions.

\textbf{Question-level clustering in significance tests.} The 1,200 poisoned triplets derive from 100 base questions, with each question contributing twelve triplets (three injection types × four noise levels). Triplets sharing a base question are not independent: they share the question text, the ground-truth answer, and the original supporting passages. We ran a question-clustered bootstrap (1,000 iterations resampling at the base-question level) for the McNemar tests reported in Section 4.4.5. The bootstrap-corrected p-values are reported alongside the uncorrected values in Table~\ref{tab:mcnemar-true-faithfulness} and discussed in the surrounding text. The directional findings hold under bootstrap correction, but the precision of the True-subset significance claims drops substantially: GPT and Gemma's True-subset effects do not survive clustered correction (bootstrap p = 0.239 and 0.259), only DeepSeek's True-subset effect approaches significance (bootstrap p = 0.099), and the most robust significant effect is the DeepSeek Survived-subset shift driven by PoisonedRAG-style triplets (bootstrap p = 0.012). We treat the bootstrap-corrected p-values as the more honest characterisation of the data.

\textbf{Filter design choices.} The statistical filter uses five engineered signals aggregated via XGBoost. The signal set is informed by the attack types in our dataset and by the multi-hop QA setting; alternative signal sets (e.g., learned representation-level features, transformer-based passage classifiers) could plausibly produce different results. The Mistral filter uses a single triplet-level prompt; alternative prompt designs could plausibly produce different recall and precision profiles.

\textbf{Single dataset version.} The experiments use a single version of the Poisoned Evaluation Dataset, with 100 base questions and 2,400 triplets. Sample-size considerations limit the precision of per-cell estimates, particularly for the smaller subsets (True n = 576 partitioned across 3 judges × 3 conditions × 2 injection types means roughly 96 triplets per granular cell). Larger datasets would tighten confidence intervals but are unlikely to change the headline directions.

\textbf{The True/Survived split is structurally asymmetric across injection types.} The decomposition cleanly separates "support overwritten" from "support intact" for replacement-based attacks (Types 1 and 2). For additive attacks (Type 3, PoisonedRAG-style), the decomposition is degenerate: all such triplets are in the Survived subset by construction. This is methodologically appropriate - additive injection by definition does not overwrite anything - but it means the True/Survived dichotomy is most informative on replacement attacks and less informative on attacks that operate by augmentation rather than replacement. Other reasonable decompositions exist (for example, partitioning by whether the judge effectively had access to a majority of supporting evidence post-filter); we report the structural definition because it is computable from labels rather than from judge behaviour, but we do not claim it is the only sensible one.

\textbf{Production deployability of the True/Survived analysis.} The decomposition into True and Survived requires per-passage poisoning labels that are available because we constructed the dataset. Real deployments do not have such labels. A practitioner cannot directly compute True-category FPR in production; they can only measure judge behaviour and infer how the filter is changing what the judge sees. This is a methodological limit on directly importing the present work's evaluation procedure into a production monitoring pipeline. What practitioners can do is apply representative test datasets at evaluation time (such as the one released with this thesis) to estimate how a candidate filter affects judge reliability before deploying.

\textbf{Justification analysis sample size.} The justification analysis is based on 28 triplets scored by a single judge under three prompt conditions. This is a qualitative supplement, not a quantitative claim, and we treat it as such throughout.

\textbf{Mechanism evidence is empirical but bounded.} Where Chapter 4 and Chapter 5 attribute filter-induced shifts to mechanisms, the principal mechanism claim (context clarification through distractor removal) is supported by direct empirical evidence: a counterfactual that refutes the collateral-damage hypothesis, a stratification that shows monotonically rising Δ with distractor-removal volume, and a length-correlation analysis showing that judges score shorter post-filter contexts as more faithful. These tests rule out the most obvious alternative explanations but do not constitute a fully causal claim. The mechanism is best characterised as the explanation most consistent with the available evidence after testing the principal alternatives. A randomised intervention varying context length while holding poison content constant would tighten this further but is outside the scope of the present work.

\section{Practical Implications}

The findings translate into several concrete recommendations for practitioners building LLM-as-a-Judge pipelines for RAG.

First, baseline judge FPRs on poisoned context are unacceptably high for security-critical deployments. Pipelines that rely on a single LLM judge to make automated quality decisions on RAG outputs are exposed to corpus-level adversarial pressure, and the exposure is not small (baseline faithfulness FPRs of 41\% to 70\% across the three judges studied). For deployments in adversarial settings - compliance monitoring, financial document processing, medical evaluation, legal review - unfiltered LLM-as-a-Judge is not a complete solution.

Second, lightweight upstream filtering does not solve the problem. The intuitive practitioner response to judge vulnerability is to filter the context before it reaches the judge. The empirical results in this thesis show that this intuition is wrong in current designs. Pre-filtering changes the residual context in ways that can make judges more confident in poisoned content, not less. A pipeline that adds pre-filtering and assumes downstream reliability has improved is making an assumption the data does not support.

Third, judge selection is a security decision. The three judges in our panel respond differently to the most sophisticated attack class. The practical implication is that pipelines should not select judges purely on cost, and should run baseline robustness measurements under representative adversarial conditions before committing to a judge for production use. The Poisoned Evaluation Dataset released with this thesis can serve as a starting point for such measurements.

Fourth, multi-judge ensembling is not a substitute for filtering. Pipelines that deploy multiple judges and rely on disagreement for detection will not see the gain they expect; the disagreement signal is far too small to support thresholded detection. Ensembling can be valuable for other reasons (calibration averaging, redundancy against API failures) but not for poisoning detection.

Fifth, judge-side hardening is a more promising direction than filter optimisation. The justification analysis suggests that explicit poison-aware prompting can substantially shift judge behaviour on PoisonedRAG-style attacks. Practitioners with pipeline access to judge prompts have a tractable mitigation that does not require building or training a filter. Whether this scales to production conditions is an empirical question, but the early signal is positive.

Sixth, evaluation methodology matters. Practitioners measuring filter effectiveness should use the True/Survived decomposition, or some operational equivalent, to avoid being misled by classifier-level F1. A filter that achieves high F1 on a held-out test split may produce post-filter residuals that judges handle worse than the unfiltered context. The right metric is the one that measures end-to-end judge reliability, not classification metrics on the filter in isolation.

---
